\documentclass[12pt,a4paper]{article}
\usepackage[polish]{babel}
\usepackage[T1]{fontenc}
\usepackage[utf8x]{inputenc}
\usepackage{hyperref}
\usepackage{url}
\usepackage[]{algorithm2e}
\usepackage{listings}

\usepackage{color}
\usepackage{listings}

\usepackage{tikz}
\usetikzlibrary{positioning,arrows.meta}

\lstloadlanguages{
  Go,        % Dodanie Go
}

\setcounter{section}{0}

\definecolor{red}{rgb}{0.6,0,0} % for strings
\definecolor{blue}{rgb}{0,0,0.6}
\definecolor{green}{rgb}{0,0.8,0}
\definecolor{cyan}{rgb}{0.0,0.6,0.6}
\lstset{
    language=Go,
    basicstyle=\footnotesize\ttfamily,
    numbers=left,
    numberstyle=\tiny,
    numbersep=5pt,
    tabsize=2,
    extendedchars=true,
    breaklines=true,
    frame=b,
    stringstyle=\color{blue}\ttfamily,
    showspaces=false,
    showtabs=false,
    xleftmargin=17pt,
    framexleftmargin=17pt,
    framexrightmargin=5pt,
    framexbottommargin=4pt,
    commentstyle=\color{green},
    morecomment=[l]{//},  % dla komentarzy liniowych
    morecomment=[s]{/*}{*/},  % dla komentarzy wieloliniowych
    showstringspaces=false,
    morekeywords={
        package, import, func, var, const, type, go, defer, return,
        if, else, switch, case, fallthrough, for, range, break, continue,
        select, chan, interface, struct, map, gofmt, len, cap, append, copy,
        error, defer, panic, recover
    },
    keywordstyle=\color{cyan},  % kolory dla słów kluczowych
    identifierstyle=\color{red},  % kolory dla identyfikatorów
    commentstyle=\color{green},  % kolory dla komentarzy
    emph={main},  % Możesz też wyświetlić główną funkcję na inny sposób
    emphstyle=\bfseries\color{magenta}  % dla wyróżnionej funkcji "main"
}

\usepackage{caption}
\DeclareCaptionFont{white}{\color{white}}
\DeclareCaptionFormat{listing}{\colorbox{blue}{\parbox{\textwidth}{\hspace{15pt}#1#2#3}}}
\captionsetup[lstlisting]{format=listing,labelfont=white,textfont=white, singlelinecheck=false, margin=0pt, font={bf,footnotesize}}


\addtolength{\hoffset}{-1.5cm}
\addtolength{\marginparwidth}{-1.5cm}
\addtolength{\textwidth}{3cm}
\addtolength{\voffset}{-1cm}
\addtolength{\textheight}{2.5cm}
\setlength{\topmargin}{0cm}
\setlength{\headheight}{0cm}

\begin{document}
	
	\title{Modelowanie i analiza systemów informatycznych\\\small{dokumentacja projektu RSP - Rozproszony system plików}}
	\author{Piotr Skowroński, Krzysztof Czuba, grupa XYZ}
	\date{\today}

	\maketitle
	\newpage
    \section{Część 1.}
\subsection{Opis programu}
Implementacja rozproszonego systemu plików, z uwzględnieniem problemu zakleszczeń i algorytmu elekcji.

System składa się z części serwerowej w której działa jeden węzeł główny (\texttt{master}) oraz wiele węzłów jedynie przechowujących dane (\texttt{storage}), oraz z web UI umożliwiającego download/upload plików.

\subsection{Wymagania wstępne i środowisko}

\subsubsection{Wymagania systemowe i narzędzia}
Do uruchomienia i wdrożenia systemu wymagane są następujące narzędzia:
\begin{itemize}
    \item \textbf{System Operacyjny:} \textbf{Linux} (testowane lokalnie; inne systemy wymagają analogicznych narzędzi).
    \item \textbf{Repozytorium:} \textbf{Git} (pobranie projektu).
    \item \textbf{Backend i węzły DFS:} \textbf{Go} (\texttt{1.25+}).
    \item \textbf{Baza Danych:} \textbf{PostgreSQL} (\texttt{5432/TCP}) -- baza danych dla API.
    \item \textbf{Frontend:} \textbf{Node.js} (LTS, np. \texttt{20.x}) i \textbf{pnpm}.
\end{itemize}

\subsubsection{Wymagania sieciowe}
Aby klaster działał poprawnie, należy zapewnić komunikację TCP:
\begin{itemize}
    \item Master powinien być osiągalny dla węzłów Storage na porcie \texttt{9000/TCP} (lub innym skonfigurowanym).
    \item Na maszynie zewnętrznej, na której działa węzeł Storage, port tego węzła (np. \texttt{9004/TCP}) musi zostać odblokowany w firewallu.
    \item Z Mastera powinno być możliwe zestawienie połączenia do \texttt{advertise-ip:port} węzła Storage.
\end{itemize}


\subsection{Instalacja i Wdrożenie (Przygotowanie)}
Sekcja ta opisuje kroki niezbędne do przygotowania środowiska i binarek systemu. Kroki te są wykonywane zazwyczaj tylko raz.

\begin{enumerate}
    \item Projekt powinien zostać pobrany:
\begin{verbatim}
git clone https://github.com/H4wk507/distributed-file-system
cd distributed-filesystem
\end{verbatim}

    \item Zależności frontendu powinny zostać zainstalowane:
\begin{verbatim}
cd frontend
pnpm install
\end{verbatim}
    
    \item PostgreSQL powinien zostać uruchomiony lokalnie na \texttt{localhost:5432} (parametry powinny być zgodne z \texttt{.env} / \texttt{backend/.env}).

    \item Migracje bazy danych powinny zostać wykonane (wymagane jest narzędzie \texttt{migrate}):
\begin{verbatim}
go install -tags 'postgres' github.com/golang-migrate/migrate/v4/cmd/migrate@latest
make migrate-up
\end{verbatim}
\end{enumerate}

\subsection{Instrukcja Uruchomienia i Obsługi}

\subsubsection{Uruchomienie lokalne (Środowisko deweloperskie)}
Po wykonaniu kroków instalacji, system w trybie deweloperskim można uruchomić za pomocą trzech terminali:
\begin{enumerate}
    \item \textbf{Uruchomienie Klastra DFS} (master + 3 storage):
\begin{verbatim}
cd backend
./scripts/start-cluster.sh
\end{verbatim}
    \item \textbf{Uruchomienie API}:
\begin{verbatim}
cd backend
go run ./cmd/api
\end{verbatim}
    \item \textbf{Uruchomienie Frontendu}:
\begin{verbatim}
cd frontend
pnpm dev
\end{verbatim}
\end{enumerate}

Po uruchomieniu:
\begin{itemize}
    \item Interfejs WWW powinien być dostępny pod adresem \texttt{http://localhost:5173}.
    \item API powinno być dostępne pod \texttt{http://localhost:8080}.
\end{itemize}

\subsubsection{Uruchomienie węzła na maszynie zewnętrznej (Storage)}

**Założenie:** Master powinien być osiągalny z maszyny zewnętrznej po TCP na porcie \texttt{9000} (lub innym skonfigurowanym).

1. Binarka węzła powinna zostać zbudowana na maszynie docelowej:
\begin{verbatim}
cd backend
go build -o bin/node ./cmd/node
\end{verbatim}

2. Węzeł storage powinien zostać uruchomiony z odpowiednimi flagami:
\begin{itemize}
    \item \texttt{-ip} było adresem nasłuchiwania (zwykle \texttt{0.0.0.0}),
    \item \texttt{-advertise-ip} było adresem maszyny zewnętrznej osiągalnym z mastera,
    \item \texttt{-seed-ip} było adresem mastera osiągalnym z maszyny zewnętrznej.
\end{itemize}

Przykład (master: \texttt{192.168.1.10:9000}, węzeł zewnętrzny: \texttt{192.168.1.20:9004}):

\begin{verbatim}
cd backend
./bin/node \
    -ip 0.0.0.0 \
    -advertise-ip 192.168.1.20 \
    -port 9004 \
    -role storage \
    -priority 4 \
    -seed-ip 192.168.1.10 \
    -seed-port 9000
\end{verbatim}

	\subsection*{Dodatkowe informacje}

    Piotr Skowroński - Architektóra systemu, praca nad testami, implementacja kluczowych algorytmów, backend, frontend.\\
    Krzysztof Czuba - backend, frontend, praca nad elementami przechowywania plików, dokumentacja.

	\newpage
	\section*{Część II}
\subsection*{Opis działania}
	Poniższa sekcja opisuje działanie systemu rozproszonego w ujęciu modelowym: zegary logiczne Lamporta, rozproszone blokady, wykrywanie zakleszczeń oraz elekcję lidera. Opis jest zgodny z implementacją w katalogu \texttt{backend/dfs/}.

	\subsubsection*{Architektura i komunikacja}
	System składa się z węzłów DFS (jeden \emph{master} oraz wiele \emph{storage}) oraz warstwy API+UI.
	\begin{itemize}
		\item Węzły DFS komunikują się po TCP poprzez wiadomości JSON (\texttt{MessageWithTime}) oraz osobny protokół binarny do strumieniowania plików.
		\item Każdy węzeł utrzymuje listę \texttt{peers} oraz okresowo wysyła \texttt{heartbeat}.
		\item Strumieniowanie plików odbywa się na porcie \texttt{port+100} (\texttt{StreamPortOffset}).
	\end{itemize}

	\subsubsection*{Zegary Lamporta}
	W systemie stosowany jest zegar logiczny Lamporta $L_i$ dla każdego węzła $i$.
	\begin{itemize}
		\item Zdarzenie lokalne / wysłanie wiadomości: $L_i := L_i + 1$.
		\item Odbiór wiadomości ze znacznikiem $t$: $L_i := \max(L_i, t) + 1$.
	\end{itemize}
	Dla uzyskania porządku całkowitego stosowane jest porównanie par $(L, id)$, gdzie \texttt{id} to UUID węzła (tie-breaker). Taki porządek jest wykorzystywany m.in. do sortowania żądań blokad.

	\subsubsection*{Rozproszone blokady (kolejki żądań + ACK)}
	Dla zasobu \texttt{resourceID} każdy węzeł utrzymuje kolejkę żądań \texttt{lockQueue[resourceID]}, sortowaną stabilnie po $(L, id)$. Węzeł może wejść do sekcji krytycznej jeśli:
	\begin{enumerate}
		\item jego żądanie jest pierwsze w kolejce,
		\item otrzymał potwierdzenia \texttt{LockAck} od wszystkich znanych \texttt{peers}.
	\end{enumerate}
	Po wejściu do sekcji krytycznej węzeł rozgłasza \texttt{LockAcquired}, a po wyjściu \texttt{LockRelease}.

	\subsubsection*{Model zakleszczeń: graf oczekiwań (wait-for graph)}
	Zakleszczenie modelowane jest grafem skierowanym $G=(V,E)$, gdzie $V$ to węzły, a krawędź $u \rightarrow v$ oznacza, że węzeł $u$ czeka na zwolnienie zasobu przez $v$. Krawędzie są dodawane, gdy:
	\begin{itemize}
		\item nadejdzie \texttt{LockRequest} dla zasobu, który ma już właściciela (\texttt{lockHolders[resourceID]}),
		\item a właściciel jest różny od węzła żądającego.
	\end{itemize}
	Master cyklicznie (co 10s) sprawdza, czy $G$ zawiera cykl. Cykl oznacza zakleszczenie.

	\subsubsection*{Rozwiązywanie zakleszczeń}
	W przypadku wykrycia cyklu wybierana jest ofiara (\emph{victim}) o najniższym priorytecie w cyklu. Master wysyła do niej komunikat \texttt{LockAbort}. Ofiara zwalnia swoje blokady (rozgłasza \texttt{LockRelease}), co usuwa krawędzie w grafie oczekiwań i przerywa cykl.

	\subsubsection*{Elekcja lidera (Bully)}
	Węzły mają priorytety (liczba całkowita). Gdy master przestaje odpowiadać (brak heartbeat), uruchamiany jest algorytm Bully:
	\begin{itemize}
		\item węzeł rozpoczynający elekcję wysyła \texttt{Election} do węzłów o wyższym priorytecie,
		\item jeśli nie ma wyższych priorytetów lub nie nadejdzie \texttt{OK} w czasie, ogłasza się liderem,
		\item lider rozgłasza \texttt{Coordinator}; węzły aktualizują role (\texttt{master}/\texttt{storage}).
	\end{itemize}
	Po elekcji nowy master uruchamia synchronizację metadanych i ewentualną naprawę replikacji.

	\subsubsection*{Rozmieszczenie i replikacja plików (consistent hashing)}
	Master utrzymuje pierścień haszujący z węzłami wirtualnymi (domyślnie 150 na węzeł). Dla identyfikatora pliku \texttt{FileID} wyliczany jest hash (CRC32) i wybierane są kolejne unikalne węzły na pierścieniu jako repliki (docelowo 3).
	Po elekcji master zbiera metadane (\texttt{MetadataRequest/Response}), buduje \texttt{globalFileIndex} oraz dla plików z niedoborem replik inicjuje kopiowanie (\texttt{ReplicateFile}).

	\subsubsection*{Przechowywanie lokalne (content-addressed storage)}
	Storage zapisuje pliki adresowane zawartością: oblicza $\mathrm{SHA256}(data)$ i zapisuje jako \texttt{data/<prefix>/<hash>.dat} wraz z metadanymi \texttt{<hash>.metadata.json}. Indeks \texttt{index.json} przechowuje mapowanie \texttt{hash -> FileMetadata}. Identyczna zawartość daje ten sam hash (deduplikacja).

	\subsubsection*{Upload/Download (streaming)}
	Upload jest dwuetapowy: (1) inicjalizacja u mastera i wybór storage nodes, (2) bezpośrednie strumieniowanie binarne do storage. Storage po zapisaniu wysyła \texttt{StreamStoreAck} do mastera. Download polega na wskazaniu przez mastera jednego z replik oraz pobraniu strumieniem binarnym po hash.

	\newpage
	
	\subsection*{Algorytmy}
Poniżej przedstawiono pseudokod kluczowych algorytmów zastosowanych w systemie.

	\subsubsection*{Aktualizacja zegara Lamporta}
	\begin{algorithm}[H]
		\KwData{Lokalny zegar $L_i$, (opcjonalnie) odebrany znacznik czasu $t$}
		\KwResult{Nowa wartość $L_i$}
		\eIf{zdarzenie lokalne lub wysłanie wiadomości}{
			$L_i \leftarrow L_i + 1$\;
		}{
			$L_i \leftarrow \max(L_i, t) + 1$\;
		}
		\caption{Zegar logiczny Lamporta.}
	\end{algorithm}

	\subsubsection*{Rozproszona blokada (kolejka Lamporta + ACK)}
	\begin{algorithm}[H]
		\KwData{Zasób \texttt{resourceID}, zbiór znanych węzłów \texttt{peers}}
		\KwResult{Wejście do sekcji krytycznej lub oczekiwanie}
				extbf{RequestLock(resourceID)}\;
		$L_i \leftarrow L_i+1$\;
		Utwórz \texttt{LockRequest(resourceID, nodeID=i, logicalTime=$L_i$)}\;
		Dodaj żądanie do \texttt{lockQueue[resourceID]} i posortuj po $(logicalTime, nodeID)$\;
		Zainicjalizuj \texttt{pendingAcks[resourceID]}\;
		Rozgłoś \texttt{LockRequest} do wszystkich \texttt{peers}\;
		\BlankLine
				extbf{OnReceiveLockRequest(req)}\;
		Dodaj \texttt{req} do lokalnej kolejki i posortuj\;
		Jeśli \texttt{lockHolders[resourceID]} istnieje i $\neq req.nodeID$, dodaj krawędź $req.nodeID \rightarrow holder$ do grafu oczekiwań\;
		Odeślij \texttt{LockAck(resourceID)} do nadawcy\;
		\BlankLine
				extbf{CanEnterCriticalSection(resourceID)}\;
		\If{moje żądanie jest pierwsze w kolejce \textbf{i} liczba ACK $\geq |peers|$}{
			\Return \texttt{true}\;
		}
		\Return \texttt{false}\;
		\BlankLine
				extbf{EnterCriticalSection(resourceID)}\;
		\If{\textbf{CanEnterCriticalSection(resourceID)} == \texttt{true}}{
			Rozgłoś \texttt{LockAcquired(resourceID)}\;
		}
		\caption{Rozproszona blokada zasobu.}
	\end{algorithm}

	\subsubsection*{Wykrywanie i usuwanie zakleszczeń (graf oczekiwań)}
	\begin{algorithm}[H]
		\KwData{Graf oczekiwań $G=(V,E)$, priorytety węzłów}
		\KwResult{Przerwanie cyklu (jeśli wykryty)}
		\If{węzeł nie jest \texttt{master}}{\Return}
		\If{\textbf{HasCycle}$(G)$ zwraca \texttt{true} oraz ścieżkę cyklu $C$}{
						exttt{victim} $\leftarrow \arg\min_{v \in C}\;priority(v)$\;
			Wyślij \texttt{LockAbort} do \texttt{victim}\;
		}
		\caption{Detekcja cyklu i wybór ofiary (master).}
	\end{algorithm}

	\subsubsection*{Elekcja lidera Bully}
	\begin{algorithm}[H]
		\KwData{Priorytet własny $p_i$, lista \texttt{peers}}
		\KwResult{Ustalenie roli \texttt{master}}
		Wyślij \texttt{Election} do wszystkich węzłów o priorytecie $> p_i$\;
		\If{nie istnieje żaden węzeł o priorytecie $> p_i$}{
			Rozgłoś \texttt{Coordinator} i oczekuj krótko na ewentualny \texttt{Coordinator} od wyższego priorytetu\;
			Ustaw rolę \texttt{master}\;
			\Return
		}
		\If{otrzymano \texttt{OK} w czasie \texttt{electionTimeout}}{
			Zakończ elekcję (wyższy priorytet przejmie rolę)\;
			\Return
		}
		Rozgłoś \texttt{Coordinator} i ustaw rolę \texttt{master}\;
		\caption{Algorytm elekcji lidera Bully.}
	\end{algorithm}

	\subsubsection*{Upload pliku (master + streaming do storage)}
	\begin{algorithm}[H]
		\KwData{\texttt{FileID}, \texttt{Filename}, \texttt{ContentType}, \texttt{Size}}
		\KwResult{Zapis pliku na $r=3$ storage nodes i aktualizacja metadanych}
			extbf{Master:}\;
		Wybierz $r$ węzłów \texttt{storage} na podstawie consistent hashing (pierścień)\;
		Utwórz \texttt{sessionID} i prześlij do wybranych węzłów \texttt{StreamUploadReady(sessionID, FileID, ...)}\;
		Zwróć klientowi listę adresów strumieniowania (\texttt{ip:(port+100)})\;
		\BlankLine
			extbf{Storage:}\;
		Po odebraniu strumienia zweryfikuj, czy \texttt{sessionID} istnieje\;
		Odczytaj dokładnie \texttt{Size} bajtów i zapisz w \texttt{LocalStorage}, licząc SHA256\;
		Wyślij do mastera \texttt{StreamStoreAck(sessionID, hash, nodeID)}\;
		\caption{Przepływ uploadu przez protokół streaming.}
	\end{algorithm}
	
	\subsection*{Bazy danych}
	System wykorzystuje relacyjną bazę danych PostgreSQL do przechowywania informacji o użytkownikach oraz metadanych plików. Sama zawartość plików przechowywana jest w klastrze DFS (węzły \texttt{storage}), natomiast baza przechowuje jedynie metadane i stan replik.

	\subsubsection*{Typy i tabele}
	Schemat jest definiowany w migracji \texttt{backend/migrations/000001\_init.up.sql}.
	\begin{itemize}
		\item \texttt{users} -- użytkownicy aplikacji:
		\begin{itemize}
			\item \texttt{id} (UUID, PK), \texttt{username} (UNIQUE), \texttt{email} (UNIQUE), \texttt{password\_hash}, \texttt{role} (ENUM \texttt{user\_role}), \texttt{created\_at}, \texttt{updated\_at}.
		\end{itemize}
		\item \texttt{files} -- metadane plików (nie dane binarne):
		\begin{itemize}
			\item \texttt{id} (UUID, PK), \texttt{name}, \texttt{size}, \texttt{hash}, \texttt{content\_type}, \texttt{owner\_id} (FK \textrightarrow{} \texttt{users.id}), \texttt{created\_at}, \texttt{updated\_at}.
		\end{itemize}
		\item \texttt{replicas} -- informacja o replikach danego pliku w klastrze DFS:
		\begin{itemize}
			\item \texttt{id} (UUID, PK), \texttt{file\_id} (FK \textrightarrow{} \texttt{files.id}), \texttt{node\_id} (UUID węzła DFS), \texttt{status} (ENUM \texttt{replica\_status}), \texttt{created\_at}, \texttt{updated\_at}.
		\end{itemize}
	\end{itemize}

	Dodatkowo zdefiniowane są typy ENUM:
	\begin{itemize}
		\item \texttt{user\_role} \;\{\texttt{user}, \texttt{admin}\},
		\item \texttt{replica\_status} \;\{\texttt{synced}, \texttt{syncing}, \texttt{failed}\}.
	\end{itemize}

	W każdej tabeli pole \texttt{updated\_at} jest automatycznie aktualizowane przez trigger \texttt{update\_updated\_at\_column()}.

	\subsubsection*{Relacje i integralność}
	\begin{itemize}
		\item Relacja \texttt{users (1) --- (N) files}: jeden użytkownik może posiadać wiele plików.
		\item Relacja \texttt{files (1) --- (N) replicas}: jeden plik może mieć wiele replik na węzłach DFS.
		\item Klucze obce mają \texttt{ON DELETE CASCADE} oraz \texttt{ON UPDATE CASCADE} (usunięcie użytkownika usuwa jego pliki; usunięcie pliku usuwa rekordy replik).
	\end{itemize}

	\subsubsection*{Diagram ERD}
	\begin{figure}[htbp]
		\centering
		\begin{tikzpicture}[
			entity/.style={draw, rounded corners, align=left, font=\small, inner sep=6pt},
			rel/.style={-Latex, thick}
		]
		\node[entity] (users) {
					extbf{users}\\
			id (PK)\\
			username (UQ)\\
			email (UQ)\\
			password\_hash\\
			role: user\_role\\
			created\_at\\
			updated\_at
		};

		\node[entity, right=3.2cm of users] (files) {
					extbf{files}\\
			id (PK)\\
			name\\
			size\\
			hash\\
			content\_type\\
			owner\_id (FK)\\
			created\_at\\
			updated\_at
		};

		\node[entity, right=3.2cm of files] (replicas) {
					extbf{replicas}\\
			id (PK)\\
			file\_id (FK)\\
			node\_id\\
			status: replica\_status\\
			created\_at\\
			updated\_at
		};

		\draw[rel] (users.east) -- node[above, font=\small]{1} node[below, font=\small]{owner\_id} (files.west);
		\draw[rel] (files.east) -- node[above, font=\small]{1} node[below, font=\small]{file\_id} (replicas.west);
		\node[font=\small] at ($(files.west)!0.5!(users.east) + (0,0.65)$) {N};
		\node[font=\small] at ($(replicas.west)!0.5!(files.east) + (0,0.65)$) {N};
		\end{tikzpicture}
		\caption{Diagram ERD bazy danych (PostgreSQL).}
	\end{figure}
    \newpage
	\subsection*{Implementacja systemu}
	Sekcja opisuje implementację systemu z perspektywy podziału na moduły, ich odpowiedzialności oraz kluczowych przepływów (DFS i API/UI).

	\subsubsection*{Struktura projektu}
	\begin{itemize}
		\item \texttt{backend/} -- implementacja DFS (węzły) oraz API (Go).
		\item \texttt{frontend/} -- interfejs WWW (React + Vite).
		\item \texttt{backend/migrations/} -- migracje bazy danych (PostgreSQL).
		\item \texttt{docker-compose.yml}, \texttt{Makefile}, \texttt{scripts/} -- uruchamianie środowiska i klastra.
	\end{itemize}

	\subsubsection*{Backend API (Go)}
	Punkt wejścia API znajduje się w \texttt{backend/cmd/api/main.go}. Warstwa API jest zorganizowana typowo dla aplikacji usługowej:
	\begin{itemize}
		\item \texttt{backend/internal/handlers/} -- kontrolery HTTP (m.in. \texttt{auth\_handler.go}, \texttt{file\_handler.go}, \texttt{metrics\_handler.go}).
		\item \texttt{backend/internal/services/} -- logika domenowa (np. autoryzacja, obsługa plików).
		\item \texttt{backend/internal/database/} -- połączenie z PostgreSQL.
		\item \texttt{backend/internal/models/} -- modele danych: \texttt{User}, \texttt{File}, \texttt{Replica}.
		\item \texttt{backend/internal/middleware/} -- middleware (np. autoryzacja JWT).
	\end{itemize}
	API pełni rolę warstwy użytkowej: rejestracja/logowanie oraz operacje na plikach (upload/download/delete) delegowane do klastra DFS.

	\subsubsection*{Węzły DFS (Go)}
	Punkt wejścia węzła DFS znajduje się w \texttt{backend/cmd/node/main.go}. Rdzeń stanowi struktura \texttt{Node} w \texttt{backend/dfs/node/node.go}, która uruchamia:
	\begin{itemize}
		\item serwer TCP do komunikacji sterującej (wiadomości JSON),
		\item serwer strumieniowania (port \texttt{port+100}) do transferu danych plików,
		\item pętle monitorujące: heartbeat, zdrowie peerów, time-outy blokad, detekcję zakleszczeń.
	\end{itemize}
\begin{itemize}
			\item \textbf{Protokół wiadomości.} Typy wiadomości i struktury danych są w \texttt{backend/dfs/common/types.go} (m.in. \texttt{MessageWithTime}, \texttt{LockRequest}). Każda wiadomość przenosi znacznik czasu Lamporta, a odbiorca aktualizuje zegar regułą $L_i := \max(L_i, t)+1$.

			\item \textbf{Elekcja lidera.} Moduł \texttt{backend/dfs/election/bully.go} implementuje Bully; po wykryciu braku heartbeat lidera uruchamiana jest elekcja, a po wygranej ustawiana rola \texttt{master} oraz wykonywana synchronizacja metadanych.

			\item \textbf{Przechowywanie danych.} \texttt{backend/dfs/storage/local\_storage.go} zapisuje pliki adresowane zawartością (SHA256) oraz metadane w plikach JSON, utrzymując indeks \texttt{index.json}. Dzięki temu identyczna zawartość mapuje się na ten sam \texttt{hash}.
			\item \textbf{Rozmieszczenie i replikacja.} \texttt{backend/dfs/hashing/hash.go} realizuje consistent hashing z węzłami wirtualnymi (domyślnie 150). Master wybiera $r=3$ węzły storage jako repliki, a po elekcji wykonuje \texttt{MetadataRequest/Response} i naprawia brakujące repliki przez \texttt{ReplicateFile}.

			\item \textbf{Zakleszczenia i blokady.} Blokady realizowane są przez kolejkę żądań sortowaną po $(L, id)$ i mechanizm ACK. Master buduje graf oczekiwań (wait-for graph) i wykrywa cykle (DFS), a następnie wysyła \texttt{LockAbort} do ofiary o najniższym priorytecie w cyklu.

			\item \textbf{Streaming.} \texttt{backend/dfs/streaming/} zawiera binarny protokół transferu (nagłówek o stałej długości + stream danych) oraz zarządzanie sesjami uploadu. Inicjalizacja sesji odbywa się u mastera, a klient wysyła dane bezpośrednio do wybranych storage nodes.	
	\end{itemize}
	\subsubsection*{Frontend (React + Vite)}
	Frontend znajduje się w \texttt{frontend/}.
	\begin{itemize}
		\item \texttt{src/pages/} -- widoki (np. logowanie, monitoring, węzły, pliki).
		\item \texttt{src/hooks/} -- logika komunikacji z API (np. pobieranie listy plików, upload/download).
		\item \texttt{src/components/} -- komponenty UI (lista plików, podgląd, upload, widok węzłów).
	\end{itemize}
	UI wykorzystuje API do operacji użytkownika, a sam transfer danych plików przebiega zgodnie z protokołem streaming (init u mastera + upload/download do storage).

	\subsubsection*{Przykładowe fragmenty kodu}
	Poniżej przedstawiono krótkie fragmenty implementacji (dla pełnego kodu patrz katalog \texttt{backend/}).
	\begin{lstlisting}
// Lamport: total order (t, nodeID)
func CompareLamport(t1 int, id1 uuid.UUID, t2 int, id2 uuid.UUID) int {
    if t1 < t2 { return -1 }
    if t1 > t2 { return 1 }
    // tie-break by node ID
    if id1.String() < id2.String() { return -1 }
    if id1.String() > id2.String() { return 1 }
    return 0
}
	\end{lstlisting}
	\subsection*{Testy}
	Testy jednostkowe zaimplementowano w Go w katalogu \texttt{backend/dfs/} (pliki \texttt{\_test.go}). Pokrywają one głównie logikę algorytmiczną (bez uruchamiania pełnego klastra).

	\subsubsection*{Uruchamianie}
	\begin{verbatim}
cd backend
go test ./...
	\end{verbatim}
	Opcjonalnie (pokrycie):
	\begin{verbatim}
go test ./... -cover
	\end{verbatim}

	\subsubsection*{Zakres}
	\begin{itemize}
		\item \texttt{backend/dfs/common/waitfor-graph\_test.go} -- testy detekcji cyklu w grafie oczekiwań (przypadki: graf pusty, pętla własna, cykle 2- i 3-węzłowe, graf rozłączny).
		\item \texttt{backend/dfs/hashing/hash\_test.go} -- testy consistent hashing (deterministyczność \texttt{HashKey}, węzły wirtualne, dodawanie/usuwanie węzłów, stabilność wyboru replik dla tego samego \texttt{FileID}).
		\item \texttt{backend/dfs/storage/local\_storage\_test.go} -- testy \texttt{LocalStorage} (\texttt{SaveFile/GetFile/DeleteFile}, idempotencja usuwania, zapis/odczyt \texttt{index.json}, deduplikacja: ta sama zawartość $\Rightarrow$ ten sam hash).
		\item \texttt{backend/dfs/node/node\_test.go} -- testy logiki węzła (kolejki/ACK blokad, integracja z grafem oczekiwań, wybór ofiary zakleszczenia, obsługa metadanych i replikacji).
	\end{itemize}

	\subsubsection*{Wynik (stan bieżący repozytorium)}
	W aktualnym stanie repozytorium polecenie \texttt{go test ./...} kończy się sukcesem (w tym pakiet \texttt{dfs/node}).
	\begin{verbatim}
ok   dfs-backend/dfs/common   (tests passed)
ok   dfs-backend/dfs/hashing  (tests passed)
ok   dfs-backend/dfs/storage  (tests passed)
ok   dfs-backend/dfs/node     (tests passed)
	\end{verbatim}

	\subsubsection*{Testy integracyjne/scenariusze (\texttt{backend/cmd/test-*})}
	W repozytorium znajdują się również programy uruchamiane jako osobne komendy w \texttt{backend/cmd/test-bully}, \texttt{backend/cmd/test-lock}, \texttt{backend/cmd/test-message}, \texttt{backend/cmd/test-sync}. Nie są one wykrywane przez \texttt{go test}, ponieważ nie mają postaci plików \texttt{\_test.go} z funkcjami \texttt{TestXxx(*testing.T)} — pełnią rolę „harnessów” do uruchomienia mini-klastra i obserwacji logów.

	Uruchamianie (w katalogu \texttt{backend/}):
	\begin{verbatim}
go run ./cmd/test-message
go run ./cmd/test-bully
go run ./cmd/test-lock
go run ./cmd/test-sync
	\end{verbatim}
	Krótko:
	\begin{itemize}
		\item \texttt{test-message} -- 2 węzły, discovery/heartbeat + obserwacja kanałów wiadomości.
		\item \texttt{test-bully} -- 3 węzły, symulacja zatrzymania lidera i przebiegu elekcji Bully.
		\item \texttt{test-lock} -- 1 węzeł i 5 węzłów, scenariusz kolejkowania żądań blokady i wejść do sekcji krytycznej.
		\item \texttt{test-sync} -- 4 węzły, „zabicie” mastera i weryfikacja synchronizacji metadanych po elekcji.
	\end{itemize}
	
	\subsection*{Eksperymenty}
	Poniżej przedstawiono propozycję eksperymentów, które pozwalają ilościowo opisać zachowanie systemu (czas reakcji, narzut algorytmów, przepustowość) oraz porównać wybrane parametry konfiguracji.

	W ramach repozytorium dostępne są gotowe komendy eksperymentalne (\texttt{backend/cmd/exp-*}), które zapisują wyniki do plików CSV w \texttt{backend/experiments/*.csv}. Na ich podstawie generowane są wykresy (\texttt{experiments/plot_csvs.py} \(\Rightarrow\) \texttt{experiments/plots/*.png}).

	\subsubsection*{Metodyka pomiaru}
	\begin{itemize}
		\item Każdy pomiar wykonywać $n\ge 10$ razy; raportować średnią, medianę oraz odchylenie standardowe.
		\item Mierzone czasy: w ms; przepustowość: $\mathrm{MB}/\mathrm{s}$; nierównomierność rozkładu: min/avg/max oraz $\sigma$.
		\item Warunki: stała liczba węzłów, stała konfiguracja sieci, brak innych obciążeń w tle.
	\end{itemize}

	\subsubsection*{Eksperyment 1: czas przejęcia lidera (Bully)}
				extbf{Cel:} zmierzyć czas od awarii mastera do ponownej dostępności usługi (wybór nowego lidera i stabilizacja).
	\begin{itemize}
		\item Narzędzie: \texttt{go run ./cmd/exp-bully}.
		\item Metryki (CSV: \texttt{bully.csv}):
		\begin{itemize}
			\item \texttt{failover\_ms} -- czas od „zabicia” mastera do momentu, gdy klaster znów odpowiada (miara dostępności; im mniej, tym lepiej).
			\item \texttt{new\_master\_priority}, \texttt{new\_master\_port} -- identyfikacja nowego lidera; służy weryfikacji poprawności Bully (czy wygrał węzeł o najwyższym priorytecie wśród żywych).
			\item \texttt{nodes} -- liczba węzłów w klastrze (parametr eksperymentu).
		\end{itemize}
		\item Zmienna: liczba węzłów (np. 3/5/7) oraz rozkład priorytetów.
	\end{itemize}
				extbf{Wyniki:} rozkład \texttt{failover\_ms} (średnia/mediana/\(\sigma\)) i wykres pudełkowy w zależności od liczby węzłów.

	\subsubsection*{Eksperyment 2: narzut mechanizmu blokad (Lamport+ACK)}
				extbf{Cel:} ocenić opóźnienie wejścia do sekcji krytycznej przy rosnącej konkurencji.
	\begin{itemize}
		\item Narzędzie: \texttt{go run ./cmd/exp-lock}.
		\item Metryki (CSV: \texttt{lock.csv}):
		\begin{itemize}
			\item \texttt{wait\_ms} -- czas od zgłoszenia żądania blokady do możliwości wejścia do sekcji krytycznej (miara narzutu blokad pod kontencją).
			\item \texttt{order} -- kolejność uzyskania blokady w danym runie (pozwala analizować fairness i potencjalne starvation).
			\item \texttt{priority}, \texttt{node\_idx} -- identyfikacja węzła i jego priorytetu (ułatwia interpretację rozkładów opóźnień).
		\end{itemize}
		\item Scenariusz: wszyscy proszą o lock na ten sam \texttt{resourceID}; warianty: 1/3/5 węzłów.
	\end{itemize}
				extbf{Wyniki:} rozkład \texttt{wait\_ms} (np. boxplot) oraz zależność opóźnień od \texttt{order}.

	\subsubsection*{Eksperyment 3: detekcja zakleszczeń (wait-for graph)}
				extbf{Cel:} zmierzyć czas przerwania zakleszczenia (czas „powrotu do pracy” po wystąpieniu cyklu oczekiwań).
	\begin{itemize}
		\item Scenariusz: co najmniej 2 węzły uzyskują lock na różne zasoby i następnie żądają locka na zasób drugiego węzła (cykl $A\rightarrow B\rightarrow A$).
		\item Narzędzie: \texttt{go run ./cmd/exp-deadlock}.
		\item Metryki (CSV: \texttt{deadlock.csv}):
		\begin{itemize}
			\item \texttt{resolve\_ms} -- czas od wywołania warunków zakleszczenia (drugie żądanie blokady) do momentu, gdy zakleszczenie zostaje przerwane i węzeł może wejść do sekcji krytycznej.
		\end{itemize}
		\item Zmienna: interwał sprawdzania grafu po stronie mastera (np. 1s/5s/10s) oraz obciążenie tłem.
	\end{itemize}
				extbf{Wyniki:} rozkład \texttt{resolve\_ms} dla różnych interwałów + komentarz kompromisu (szybkość reakcji vs narzut okresowej analizy grafu).

	\subsubsection*{Eksperyment 4: przepustowość upload/download (streaming)}
				extbf{Cel:} ocenić przepustowość protokołu streaming dla różnych rozmiarów plików.
	\begin{itemize}
		\item Narzędzie: \texttt{go run ./cmd/exp-streaming}.
		\item Scenariusz: upload i download plików o rozmiarach np. 1MB/10MB/100MB.
		\item Metryki (CSV: \texttt{streaming.csv}):
		\begin{itemize}
			\item \texttt{size\_bytes} -- rozmiar pliku w bajtach (parametr eksperymentu).
			\item \texttt{upload\_ms}, \texttt{upload\_mbps} -- czas i efektywna przepustowość uploadu (miara kosztu replikacji i narzutu protokołu).
			\item \texttt{download\_ms}, \texttt{download\_mbps} -- czas i przepustowość downloadu (miara ścieżki odczytu).
		\end{itemize}
	\end{itemize}
				extbf{Wyniki:} wykres \texttt{upload\_mbps}/\texttt{download\_mbps} w funkcji \texttt{size\_bytes} oraz porównanie wariantów liczby replik.

	\subsubsection*{Szablon tabeli wyników}
	\begin{center}
	\begin{tabular}{|l|c|c|c|c|}
	\hline
	Eksperyment & Parametr & Średnia & Mediana & $\sigma$ \\
	\hline
	Bully failover & $N=\dots$ & & & \\
	Lock latency & $N=\dots$ & & & \\
	Deadlock detect & $\Delta t=\dots$ & & & \\
	Streaming & size=\dots & & & \\
	\hline
	\end{tabular}
	\end{center}
	\newpage
	\section*{Pełen kod aplikacji}
\begin{lstlisting}
Tutaj wklejamy pelen kod. 
\end{lstlisting}
\end{document}